\documentclass[11pt]{article}

% ---------------------------------------------------------------------------
% Packages
% ---------------------------------------------------------------------------
\usepackage[margin=1in]{geometry}
\usepackage{amsmath, amssymb}
\usepackage{graphicx}
\usepackage{booktabs}
\usepackage{array}
\usepackage{float}
\usepackage[hidelinks]{hyperref}

% Figures live inside the three experiment folders (note the spaces in the names).
% graphicspath entries are brace-delimited, so the spaces are handled correctly.
\graphicspath{%
  {Gaussian testing/results/figures/}%
  {HighDim testing/results/figures/}%
  {Advanced processes testing/results/figures/}%
}

\newcommand{\code}[1]{\texttt{#1}}
\newcommand{\dd}{\mathrm{d}}
\newcommand{\E}{\mathbb{E}}
\newcommand{\R}{\mathbb{R}}
\newcommand{\Var}{\operatorname{Var}}
\newcommand{\Cov}{\operatorname{Cov}}
\newcommand{\sig}{\mathbb{S}}

\title{\textbf{Characterising Stochastic Processes with Path Signatures}\\[4pt]
\large From One-Dimensional Moment Recovery to High-Dimensional Path Geometry\\
and Advanced Financial Processes}
\author{PathLaw project}
\date{\today}

\begin{document}
\maketitle

\begin{abstract}
This report documents a staged study of how \emph{truncated path signatures}
characterise stochastic processes, together with the theory and implementation
behind it. On the theory side we set out the stochastic differential equations
(SDEs) and basic properties of the models we study --- geometric Brownian motion,
Heston stochastic volatility, the Merton and Bates jump-diffusions, fractional
Brownian motion and rough Bergomi volatility --- and give a self-contained account
of path signatures, their truncation, and the precise sense in which the
\emph{expected signature} encodes the moments and law of a process. On the
implementation side we detail how the two signature libraries (\code{pysiglib} and
\code{roughpy}) were cross-validated, how the batched, memory-flat Monte-Carlo
engine simulates each process, and how moments and signatures are estimated and
compared. On the results side we confirm, on one-dimensional Gaussian and
non-Gaussian processes, that each moment becomes visible at its own signature
level (mean~$\to$~1, variance~$\to$~2, skewness~$\to$~3, kurtosis~$\to$~4); we
recover cross-channel covariance geometry (level~2+) in high dimension and
benchmark it against classical estimators; and we show that on advanced financial
processes the signature separates stochastic volatility, jumps and, most
distinctively, \emph{path memory} (long memory and roughness) that a moment or
covariance table cannot see. All signature summaries are streaming and verified
flat at $\approx 0.46$~GB for $100{,}000$ paths.
\end{abstract}

\tableofcontents

% ===========================================================================
\section{Introduction}
% ===========================================================================
The guiding research question, posed by the project supervisor, is:
\begin{quote}
\emph{``We want to check the way signatures characterise processes by checking
how they recover their moments. Note, you have to truncate the level of the
signatures. Start with a Gaussian process (only two moments) and complexify the
processes (higher moments) and test the truncation levels.''}
\end{quote}
In plain terms: \textbf{can truncated path signatures characterise a stochastic
process by recovering its moment information (mean, variance, \ldots), its
geometry (covariance, co-movement) in higher dimensions, and its path memory once
we leave the Gaussian/Markov world?}

\paragraph{Four departures from the Black--Scholes benchmark.}
The classical benchmark for an asset price is geometric Brownian motion
(Black--Scholes): constant volatility and independent Gaussian log-returns. Real
markets violate this in four qualitatively different ways, and these are the
``cases'' this report studies:
\begin{enumerate}
  \item \textbf{Stochastic volatility.} Volatility is itself random and mean
        reverts, producing volatility clustering and (with leverage) skew. The
        canonical model is \textbf{Heston}.
  \item \textbf{Jumps (and jump-diffusion).} Prices move discontinuously,
        creating skew and heavy tails on top of a diffusion. The canonical models
        are \textbf{Merton} (jumps on a constant-vol diffusion) and \textbf{Bates}
        (jumps on Heston stochastic vol --- the combined jump-diffusion).
  \item \textbf{Non-Markovian dynamics / long memory.} Increments carry
        autocorrelation, so the future depends on the whole past. The canonical
        driver is \textbf{fractional Brownian motion}.
  \item \textbf{Rough volatility.} Volatility is not merely stochastic but
        \emph{rough} --- its trajectories are far less regular than Brownian
        motion and highly autocorrelated. The canonical model is
        \textbf{rough Bergomi}.
\end{enumerate}
Part~I (Theory) writes down the SDE and basic properties of each, then develops
signatures and how they access moments. Part~II (Implementation) details the
library validation, the Monte-Carlo engine and the estimation/comparison
machinery. Part~III (Results) presents the experiments, staged from simple to
advanced:
\begin{enumerate}
  \item \textbf{Stage~1 (1D Gaussian).} Recover the mean and variance.
  \item \textbf{Stage~2 (1D non-Gaussian).} Recover skewness and kurtosis with
        matched mean and variance.
  \item \textbf{Stage~3 (high-dimensional Gaussian).} Recover cross-channel
        covariance geometry and benchmark against classical estimators.
  \item \textbf{Stage~4 (advanced processes).} Separate the four departures above
        and locate each feature by signature level.
\end{enumerate}
All signatures are computed with \code{pysiglib}; \code{roughpy} is used as an
independent cross-check. Every experiment is synthetic, on data whose true
moments and geometry are known, so the signatures can be checked against a known
answer.

% ###########################################################################
\part{Theory}
% ###########################################################################

% ===========================================================================
\section{Stochastic models and their SDEs}\label{sec:models}
% ===========================================================================
We work throughout with the \emph{log-price} $X_t = \log S_t$ on a filtered
probability space $(\Omega,\mathcal{F},(\mathcal{F}_t),\mathbb{P})$, with
$W,\,Z,\,W^S,\,W^v$ standard Brownian motions and $N$ a Poisson process. All
parameters are annualised and the horizon is $T$ years.

\subsection{The Black--Scholes benchmark (GBM)}
Geometric Brownian motion is
\begin{equation}
  \dd S_t = \mu\, S_t\,\dd t + \sigma\, S_t\,\dd W_t,
  \qquad\Longleftrightarrow\qquad
  \dd X_t = \Big(\mu - \tfrac12\sigma^2\Big)\dd t + \sigma\,\dd W_t .
\end{equation}
\textbf{Properties.} $X$ has \emph{independent, stationary, Gaussian} increments:
$X_{t+\Delta}-X_t \sim \mathcal{N}\big((\mu-\tfrac12\sigma^2)\Delta,\ \sigma^2\Delta\big)$.
It is a Markov process, a (time-changed) L\'evy process and, after discounting, a
martingale. Its law is fixed by just two moments; skewness and excess kurtosis are
zero. This is the null model against which every richer process is compared, and
$\Var(X_T-X_0)=\sigma^2 T$ is the yardstick we match the other processes to.

\subsection{Stochastic volatility: the Heston model}
Heston~\cite{Heston1993} makes the variance $v_t$ a mean-reverting square-root
(Cox--Ingersoll--Ross~\cite{CIR1985}) process correlated with the price:
\begin{align}
  \dd X_t &= \Big(\mu - \tfrac12 v_t\Big)\dd t + \sqrt{v_t}\,\dd W_t^{S}, \\
  \dd v_t &= \kappa(\theta - v_t)\,\dd t + \xi\sqrt{v_t}\,\dd W_t^{v},
  \qquad \dd\langle W^{S}, W^{v}\rangle_t = \rho\,\dd t .
\end{align}
Here $\kappa$ is the mean-reversion speed, $\theta$ the long-run variance, $\xi$
the vol-of-vol and $\rho$ the leverage correlation. \textbf{Properties.} The pair
$(X,v)$ is jointly Markov, but $X$ \emph{alone} is not --- it depends on the latent
$v$. The variance stays positive when the \emph{Feller condition}
$2\kappa\theta \ge \xi^2$ holds (here $2\cdot2\cdot0.04 = 0.16 \ge 0.09 = 0.3^2$),
and its stationary law is Gamma. Stochastic volatility produces \emph{volatility
clustering} (autocorrelated squared returns) and, through $\rho<0$, negative
skew and semi-heavy tails --- excess kurtosis without the true jumps of the next
model. Increments are \emph{not} independent, so $X$ is not a L\'evy process.

\subsection{Jumps: the Merton jump-diffusion}
Merton~\cite{Merton1976} adds a compound-Poisson jump term to a constant-vol
diffusion. With jump times of intensity $\lambda$ and i.i.d.\ log jump sizes
$J_i\sim\mathcal{N}(m,\delta^2)$,
\begin{equation}
  X_t = X_0 + \Big(\mu - \tfrac12\sigma^2 - \lambda\bar\kappa\Big)t
        + \sigma W_t + \sum_{i=1}^{N_t} J_i,
  \qquad \bar\kappa = \E\!\big[e^{J}-1\big] = e^{\,m+\frac12\delta^2}-1,
\end{equation}
where $-\lambda\bar\kappa$ is the martingale compensator that removes the drift the
jumps would otherwise inject. \textbf{Properties.} $X$ is a \emph{L\'evy process}:
increments are independent and stationary, and it is Markov. Over a step $\Delta$
the increment has mean $(\mu-\tfrac12\sigma^2-\lambda\bar\kappa)\Delta+\lambda m\Delta$,
variance $\sigma^2\Delta+\lambda(m^2+\delta^2)\Delta$, \emph{negative skew} when
$m<0$, and \emph{positive excess kurtosis} $\propto \lambda(3\delta^4+6m^2\delta^2+m^4)$
--- genuinely heavy (leptokurtic) tails absent from any diffusion.

\subsection{Stochastic volatility with jumps: the Bates model}
Bates~\cite{Bates1996} combines the two, superimposing Merton jumps on Heston
stochastic volatility:
\begin{equation}
  \dd X_t = \Big(\mu - \tfrac12 v_t - \lambda\bar\kappa\Big)\dd t
            + \sqrt{v_t}\,\dd W_t^{S} + \dd\!\Big(\textstyle\sum_{i=1}^{N_t} J_i\Big),
  \qquad
  \dd v_t = \kappa(\theta - v_t)\,\dd t + \xi\sqrt{v_t}\,\dd W_t^{v}.
\end{equation}
\textbf{Properties.} It inherits volatility clustering and leverage from Heston and
skew/heavy tails from Merton: the jumps dominate short-horizon tail behaviour while
stochastic volatility shapes the term structure of higher moments. It is the
richest of the diffusion-based models here and neither Markov in $X$ alone nor
L\'evy.

\subsection{Non-Markovian dynamics: fractional Brownian motion}
Fractional Brownian motion (fBM)~\cite{MandelbrotVanNess1968} $B^{H}$ with Hurst
index $H\in(0,1)$ is the centred Gaussian process with $B^H_0=0$ and
\begin{equation}
  \E\!\big[B^{H}_t B^{H}_s\big]
  = \tfrac12\big(|t|^{2H} + |s|^{2H} - |t-s|^{2H}\big).
\end{equation}
We take the log-price to be $X_t = c\,B^{H}_t$ (with $c$ chosen so
$\Var(X_T)=\sigma^2 T$, matching Black--Scholes). \textbf{Properties.} $B^H$ is
\emph{self-similar}, $B^{H}_{at}\stackrel{d}{=}a^{H}B^{H}_t$, and has
\emph{stationary} increments. The increment process (fractional Gaussian noise,
fGn) has autocovariance
\begin{equation}\label{eq:fgn}
  \gamma(k) = \tfrac12\big(|k+1|^{2H} - 2|k|^{2H} + |k-1|^{2H}\big),
\end{equation}
so for $H>\tfrac12$ increments are \emph{positively correlated with long-range
dependence} ($\sum_k\gamma(k)=\infty$, persistent/trending), for $H<\tfrac12$ they
are anti-persistent, and $H=\tfrac12$ recovers Brownian motion. Crucially, for
$H\ne\tfrac12$ fBM is \emph{neither Markov nor a semimartingale}: the future
depends on the entire past. Its \emph{marginal} law is still Gaussian, so moments
alone cannot distinguish it from GBM once the variance is matched --- only the path
memory differs. We use $H=0.70$.

\subsection{Rough volatility: the rough Bergomi model}
Empirically, log-volatility behaves like fBM with a \emph{small} Hurst index
$H\approx0.1$ --- ``volatility is rough''~\cite{GatheralJaissonRosenbaum2018}. The
rough Bergomi model~\cite{BayerFrizGatheral2016} drives the instantaneous variance
by a rough Volterra (Riemann--Liouville) process
$\widetilde W^{H}_t = \sqrt{2H}\int_0^t (t-s)^{H-\frac12}\,\dd Z_s$:
\begin{align}
  v_t &= \xi_0\,\exp\!\Big(\eta\,\widetilde W^{H}_t - \tfrac12\eta^2 t^{2H}\Big), \\
  \dd X_t &= -\tfrac12 v_t\,\dd t
             + \sqrt{v_t}\,\big(\rho\,\dd Z_t + \sqrt{1-\rho^2}\,\dd Z_t^{\perp}\big),
\end{align}
with forward variance $\xi_0$, vol-of-vol $\eta$, roughness $H\in(0,\tfrac12)$ and
leverage $\rho$. The term $-\tfrac12\eta^2 t^{2H}$ normalises $\E[v_t]=\xi_0$.
\textbf{Properties.} Volatility paths have H\"older regularity $H-\varepsilon$ ---
far rougher than Brownian --- and are strongly (anti-)autocorrelated at short lags.
The variance is non-Markov and non-semimartingale; the model reproduces the steep
short-maturity skew that classical stochastic-volatility models miss. We use
$H=0.10$, $\eta=0.8$.

\subsection{Summary of properties}
\begin{table}[H]
\centering
\caption{Qualitative properties of the six models. ``Markov'' refers to the
observable log-price $X$ alone; jump/rough/fractional models add structure that a
two-moment description cannot capture.}
\renewcommand{\arraystretch}{1.15}
\begin{tabular}{@{}lccccc@{}}
\toprule
Model & Markov in $X$ & Increment law & Memory & Skew & Heavy tails \\
\midrule
Black--Scholes & yes & Gaussian, i.i.d.\ & none & no & no \\
Heston         & no (needs $v$) & non-Gaussian & clustering & yes ($\rho$) & mild \\
Merton         & yes (L\'evy) & Gaussian $+$ jumps & none & yes ($m<0$) & yes \\
Bates          & no & jumps $+$ SV & clustering & yes & yes \\
fBM ($H{=}0.7$)& no & Gaussian, correlated & long memory & no & no \\
rough Bergomi  & no & non-Gaussian & rough/short memory & yes ($\rho$) & yes \\
\bottomrule
\end{tabular}
\end{table}

% ===========================================================================
\section{Path signatures}\label{sec:sig}
% ===========================================================================
A \emph{path} is a continuous map $x:[0,T]\to\R^d$; discretely, a sequence of
vectors. The \emph{signature}~\cite{Chen1957,Lyons1998,LyonsCaruanaLevy2007} is the
collection of all iterated integrals of the path. For a multi-index
$(i_1,\dots,i_k)\in\{1,\dots,d\}^k$ the corresponding coordinate is
\begin{equation}
  S(x)^{\,i_1\cdots i_k}
  = \idotsint\limits_{0<t_1<\cdots<t_k<T}
    \dd x^{i_1}_{t_1}\,\dd x^{i_2}_{t_2}\cdots \dd x^{i_k}_{t_k}\ \in\ \R,
\end{equation}
and the level-$k$ tensor collects all $d^{k}$ of them,
\begin{equation}
  S(x)^{(k)} = \idotsint\limits_{0<t_1<\cdots<t_k<T}
  \dd x_{t_1}\otimes\cdots\otimes\dd x_{t_k}\ \in\ (\R^{d})^{\otimes k}.
\end{equation}
The full signature $S(x) = \big(1, S^{(1)}, S^{(2)}, \dots\big)$ lives in the tensor
algebra $T((\R^d))$ and begins with the empty-word constant $1$.

\paragraph{Two distinct axes.}
It is essential not to confuse two quantities:
\begin{center}
\begin{tabular}{@{}lll@{}}
\toprule
Concept & Meaning & Range in this study \\
\midrule
Path dimension $d$ & number of channels / assets / factors & $1$ (or $2$ with time), then $20$/$10$ \\
Signature level $k$ & order of the iterated-integral feature & up to $3$ or $5$ \\
\bottomrule
\end{tabular}
\end{center}
At level $k$ in dimension $d$ there are exactly $d^{k}$ coordinates, so the
signature truncated at depth $N$ has length $d + d^2 + \cdots + d^{N}$.

\paragraph{Algebraic structure.}
Two facts drive everything below. (i) \emph{Chen's identity}: signatures are
multiplicative under path concatenation, $S(x\ast y)=S(x)\otimes S(y)$, which is
what makes them computable incrementally and streaming. (ii) The \emph{shuffle
identity}: the pointwise product of two signature coordinates is a linear
combination (the shuffle) of higher signature coordinates. Consequently
\emph{every polynomial in the increments of the path is a linear functional of the
signature}. This is the algebraic reason signatures give access to moments
(Section~\ref{sec:moments}).

\subsection{Truncation and its implications}
The signature is infinite, so in practice we \textbf{truncate} at a finite depth
$N$, keeping levels $1,\dots,N$. Two properties make this benign and quantifiable:
\begin{itemize}
  \item \textbf{Factorial decay.} For a path of length (1-variation) $L$,
        $\big\|S(x)^{(k)}\big\| \le L^{k}/k!$~\cite{LyonsCaruanaLevy2007}. Higher
        levels shrink factorially, so a modest $N$ already captures most of the
        analytic information --- truncation error is controlled a priori.
  \item \textbf{Uniqueness.} The (full) signature determines a path up to
        tree-like equivalence and reparametrisation~\cite{HamblyLyons2010,
        BoedihardjoGengLyonsYang2016}; adding a monotone \emph{time channel}
        removes the reparametrisation ambiguity, so a time-augmented path is
        recovered from its signature. This is why we time-augment the
        one-dimensional processes in Stage~4.
\end{itemize}
The \emph{implications} of truncation are the theme of this project: a depth-$N$
signature can only resolve structure up to order $N$. A Gaussian needs $N=2$; a
skewed process needs $N\ge3$; a heavy-tailed one needs $N\ge4$. Truncating too
shallow makes two processes that differ only in a higher moment look identical ---
exactly the phenomenon the experiments probe. (For dimensionality one may instead
use the \emph{log-signature}, a compressed basis of the free Lie algebra whose
size is given by the Witt formula; we use the full signature throughout for
simplicity, since our effective dimensions are small.)

\subsection{The expected signature and characterisation of laws}
For a \emph{random} path $X$ we form the \textbf{expected signature}
$\Phi(X) = \E[S(X)]$, the average of the signature over the law of the path. Under
moment-type conditions the expected signature \emph{determines the law of the
process}~\cite{Fawcett2003,LevinLyonsNi2013,ChevyrevOberhauser2022}: it is the
path-space analogue of the classical moment problem, where the sequence of moments
pins down a distribution. This is the precise sense in which ``signatures
characterise a process''. Comparing two processes therefore reduces to comparing
their expected signatures, and truncating at level $N$ compares them up to
$N$-th-order information --- which is what we do, level by level, throughout the
report. The signature method and these characterisation results are surveyed in
\cite{ChevyrevKormilitzin2016}; efficient differentiable implementations and
kernels are developed in \cite{KidgerBonnier2019,KidgerLyons2021,Salvi2021}.

% ===========================================================================
\section{Accessing moments with signatures}\label{sec:moments}
% ===========================================================================
\subsection{From signatures to moments (mathematically)}
\paragraph{The one-dimensional case.}
For a scalar path $X$ the level-$k$ signature telescopes:
\begin{equation}\label{eq:1d}
  S(X)^{(k)}
  = \idotsint\limits_{0<t_1<\cdots<t_k<T}\dd X_{t_1}\cdots\dd X_{t_k}
  = \frac{\big(X_T - X_0\big)^{k}}{k!}.
\end{equation}
Taking expectations,
$\Phi(X)^{(k)} = \E\!\big[(X_T-X_0)^k\big]/k!$, so the level-$k$ coordinate of the
expected signature is exactly the $k$-th \emph{raw moment} of the total increment,
divided by $k!$. Hence:
\begin{center}
level $1\leftrightarrow$ mean,\quad level $2\leftrightarrow$ 2nd moment,\quad
level $3\leftrightarrow$ 3rd moment (skew),\quad level $4\leftrightarrow$ 4th
moment (kurtosis).
\end{center}
This is why, in 1D, comparing expected signatures level by level recovers moment
differences, and why higher 1D levels carry little \emph{new} information (they are
powers of a single number).

\paragraph{Time augmentation and path dependence.}
Augmenting to the two-dimensional path $t\mapsto(t, X_t)$ adds \emph{cross} terms
such as $\int_0^T t\,\dd X_t$ and $\int_0^T X_t\,\dd t$, which depend on the
\emph{order} in which the increments occur, not merely on the endpoint. These are
what let the signature separate processes with identical endpoint moments but
different memory --- fractional and rough processes in Stage~4.

\paragraph{The multi-dimensional case.}
For $d>1$ the level-2 term $\int\dd X^{i}\,\dd X^{j}$ splits into a symmetric part
$\tfrac12\,\Delta X^{i}\Delta X^{j}$ (whose expectation gives the
\emph{covariance}) and an antisymmetric part, the \emph{L\'evy area}
$\tfrac12(\int X^i\dd X^j - \int X^j\dd X^i)$, which encodes ordered co-movement a
symmetric covariance matrix cannot. More generally, by the shuffle identity every
mixed moment of the increments is a linear functional of the expected signature, so
the truncated expected signature is a rich, ordered generalisation of the moment
vector.

\subsection{Monte-Carlo estimation of the process moments}
Because the processes are simulated, their moments are estimated by Monte Carlo.
From $n$ i.i.d.\ paths we form the endpoint increments $\{\Delta X^{(p)}\}_{p=1}^n$
and estimate raw/central moments by sample averages: the sample mean
$\hat\mu=\frac1n\sum_p\Delta X^{(p)}$, the sample variance, and standardised
central moments for skewness and excess kurtosis. In the multivariate case we
estimate the mean vector $\hat\mu=\frac1n\sum_p x^{(p)}$ and covariance
$\widehat\Sigma=\frac1n\sum_p x^{(p)}x^{(p)\top}-\hat\mu\hat\mu^{\top}$. These are
the ``textbook'' summaries a quant would reach for; their sampling error scales
like $n^{-1/2}$, and \emph{higher} moments are noisier, which is why Stage~2 uses
more paths (a running sum of many increments is pushed toward Gaussianity by the
central limit theorem, diluting the higher-moment signal).

\subsection{Comparing moments and signatures}\label{sec:measures}
We compare processes under two lenses and report several normalisations.
\begin{itemize}
  \item \textbf{Statistical distances.} $L^2$ distance between mean vectors and
        Frobenius distance between covariance matrices,
        $\|\hat\mu_A-\hat\mu_B\|_2$ and
        $\|\widehat\Sigma_A-\widehat\Sigma_B\|_F$.
  \item \textbf{Signature distances.} $L^2$ distance between expected signatures.
        Because a level-$k$ block has $d^{k}$ coordinates that are also
        numerically larger, a raw distance conflates ``real difference'' with
        ``scale''. We therefore report, per level: (1) the \textbf{raw} $L^2$
        distance; (2) the \textbf{level-normalised} distance, dividing by
        $\sqrt{d^{k}}$ (a per-coordinate RMS that makes levels of different width
        comparable); and (3) the \textbf{standardised} distance, dividing each
        coordinate by the baseline process's signature-coordinate standard
        deviation before measuring, which removes the magnitude effect and
        exposes the true switch-on structure. The standardised view is the one we
        plot.
\end{itemize}

% ###########################################################################
\part{Implementation}
% ###########################################################################

% ===========================================================================
\section{Validation of the signature libraries}\label{sec:validation}
% ===========================================================================
Signatures are easy to get subtly wrong: the truncation convention, the leading
constant term, the channel ordering and the time parametrisation are all places a
pipeline can silently disagree with the mathematics. We therefore cross-validate
the two libraries before trusting any number.

\paragraph{Settings.} Signatures are computed with \code{pysiglib} via
\code{pysiglib.sig(paths, degree=$N$, time\_aug, end\_time, n\_jobs=-1)}. The
convention is \code{scalar\_term=False}, i.e.\ the returned vector omits the leading
$1$ and is laid out as the concatenation
$[\,\text{level }1\,|\,\text{level }2\,|\cdots|\,\text{level }N\,]$ with level $k$
occupying $d^{k}$ consecutive coordinates; $\code{n\_jobs}=-1$ uses all cores.
Arrays are passed as C-contiguous \code{float64} (else the library warns and
copies). Time augmentation, when used, appends a single monotone channel on
$[\,\text{end\_time}/T,\ \text{end\_time}\,]$.

\paragraph{Cross-check.} We recompute the Stage-1 expected signatures at every
level with both \code{pysiglib} and \code{roughpy}, matching all conventions (no
scalar term, identical channel order, identical time parametrisation and depth).
The two libraries agree to numerical precision: across all processes and levels the
largest disagreement is $\approx 2\times10^{-13}$, i.e.\ floating-point rounding.
This confirms the truncation layout and parametrisation before we scale up, and is
the reason the level-indexing used by every distance table (Section~\ref{sec:measures})
can be trusted.

% ===========================================================================
\section{The Monte-Carlo simulation engine}\label{sec:mc}
% ===========================================================================
\subsection{Streaming, batched design}
No experiment ever stores all paths or all signatures at once. For each process we
loop: (i) simulate a batch of paths (and their increments); (ii) update the
classical estimators; (iii) compute the batch's signatures and fold them into
fixed-length running sums; (iv) discard the batch. Memory is therefore independent
of the number of paths --- verified flat at $\approx0.46$~GB for a $100{,}000$-path
Stage-3 run. Reproducibility is guaranteed by seeding a dedicated
\code{numpy.random.default\_rng} per (configuration, process) with deterministic
offsets (Python's built-in \code{hash} is randomised per process and is
deliberately \emph{not} used).

\subsection{Discretisation schemes}
Each model uses an explicit, documented scheme on a grid of
\code{time\_steps} points with step $\dd t = T/\code{time\_steps}$:
\begin{itemize}
  \item \textbf{GBM / Merton diffusion:} exact Gaussian increments
        $(\mu-\tfrac12\sigma^2)\dd t + \sigma\sqrt{\dd t}\,Z$.
  \item \textbf{Heston / Bates variance (CIR):} Euler \emph{full-truncation}
        scheme~\cite{LordKoekkoekVanDijk2010,Andersen2008}
        $v_{t+1}=v_t+\kappa(\theta-v_t^{+})\dd t+\xi\sqrt{v_t^{+}}\sqrt{\dd t}\,Z_v$
        with the variance floored at $0$ wherever it is used, and the correlated
        price innovation $Z_s=\rho Z_v+\sqrt{1-\rho^2}Z_\perp$.
  \item \textbf{Merton / Bates jumps:} per step draw $N_t\sim\text{Poisson}(\lambda\dd t)$;
        the total log jump is $N_t\,m+\sqrt{N_t}\,\delta\,Z$ (a sum of $N_t$
        i.i.d.\ normals), less the compensator $\lambda\bar\kappa\,\dd t$.
  \item \textbf{fBM:} exact fractional Gaussian noise via the Cholesky factor $L$
        of the Toeplitz covariance built from~\eqref{eq:fgn}; a batch is
        $Z L^\top$. The factor depends only on $(\code{time\_steps},H)$ and is
        cached across batches. The scale $\sigma\sqrt{T}/\code{time\_steps}^{H}$
        matches the terminal variance to Black--Scholes.
  \item \textbf{Rough Bergomi:} the fractional process is built on the grid from
        the same Cholesky factor ($H=0.1$); we form $v_t$ by the martingale-
        normalised exponential and correlate the spot Brownian with the per-step
        vol innovation (leverage). Building the fractional process exactly in
        distribution while correlating the spot with the step innovation is a
        deliberate simplification of the full Volterra construction, adequate for
        this synthetic testbed.
\end{itemize}
The 1-D log-price is the cumulative sum of the increments; a monotone time channel
is appended to the \emph{path} (not the increments) when time augmentation is on.

\subsection{Settings and parameters}
\begin{table}[H]
\centering
\caption{Global Monte-Carlo settings per stage (seed $42$ throughout).}
\begin{tabular}{@{}lccccc@{}}
\toprule
Stage & Paths & Steps & Horizon & Depth & Dimension \\
\midrule
1 (1D Gaussian) & $1{,}000$ & $20$ & --- & $\le5$ & $1$ \\
2 (1D moments)  & $4{,}000$ & $20$ & --- & $\le5$ & $1$ \\
3 (high-dim)    & $10{,}000$ & $300$ & --- & $3$ / $5$ & $20$ / $10$ \\
4 (advanced)    & $5{,}000$ & $100$ & $1$ yr & $5$ & $1$ ($2$ with time) \\
\bottomrule
\end{tabular}
\end{table}
\begin{table}[H]
\centering
\caption{Stage-4 model parameters (annualised; baseline volatility $\sigma=0.2$,
so $\sigma^2=0.04$). Batch size $500$; time augmentation on.}
\renewcommand{\arraystretch}{1.15}
\begin{tabular}{@{}ll@{}}
\toprule
Model & Parameters \\
\midrule
Black--Scholes & $\sigma=0.2,\ \mu=0$ \\
Heston         & $v_0=\theta=0.04,\ \kappa=2.0,\ \xi=0.3,\ \rho=-0.7$ \\
Merton         & $\sigma=0.2,\ \lambda=1.0,\ m=-0.10,\ \delta=0.15$ \\
Bates          & Heston $+$ Merton parameters above \\
fBM            & $H=0.70$, scaled to $\Var(X_T)=\sigma^2 T$ \\
rough Bergomi  & $H=0.10,\ \xi_0=0.04,\ \eta=0.8,\ \rho=-0.7$ \\
\bottomrule
\end{tabular}
\end{table}

% ===========================================================================
\section{Computing and comparing moments with the libraries}\label{sec:compute}
% ===========================================================================
\subsection{Streaming expected signature}
For each process a \code{RunningSignatureStats} accumulator holds three
fixed-length quantities: the running sum of signature coordinates, the running sum
of squares, and the path count. From these we recover, without ever materialising
all signatures, the \textbf{expected signature} (sum\,/\,count) and the
\textbf{per-coordinate standard deviation} (from
$\Var = \E[x^2]-\E[x]^2$, clipped at $0$), the latter being the baseline scale used
for the standardised distances. At $d=10$, depth~5 a single signature has
$111{,}110$ numbers ($\approx0.9$~MB); storing all of them for $10{,}000$ paths
would need $\approx8.9$~GB, so this streaming is what makes the experiment feasible.

\subsection{Streaming statistical estimators}
The classical moments are estimated the same streaming way: running sums
$\sum x$ and $\sum x x^{\top}$ (length $d$ and $d\times d$) give the mean and
covariance at the end via $\widehat\Sigma=\frac1n\sum xx^\top-\hat\mu\hat\mu^\top$.
Terminal statistics use one sample per path (the endpoint); increment statistics
use every $(\text{path},\text{step})$ pair.

\subsection{Distance measures in practice}
Given the per-process expected signatures and statistical summaries, the runner
builds three tables implementing Section~\ref{sec:measures}: a statistical-distance
table ($L^2$ on means, Frobenius on covariances), a whole-signature distance table
(raw and level-normalised), and a level-wise table (raw, level-normalised and
standardised by the baseline coordinate std). Because levels occupy disjoint
coordinate blocks, the whole-signature distance truncated at depth $k$ equals
$\sqrt{\sum_{j\le k}(\text{raw level-}j\text{ distance})^2}$, so a single level
table yields the distance-versus-depth curves without recomputation.

% ###########################################################################
\part{Results}
% ###########################################################################

% ===========================================================================
\section{Stage 1: one-dimensional Gaussian processes}
% ===========================================================================
We simulate three Gaussian random walks (1000 paths, 20 time steps, seed 42),
changing one property at a time so each comparison isolates one moment.

\begin{table}[H]
\centering
\caption{Stage 1 processes and the empirical moments of their increments. The
simulator reproduces the intended mean and variance.}
\begin{tabular}{@{}lcc cc@{}}
\toprule
Process & True mean & True variance & Emp.\ mean & Emp.\ variance \\
\midrule
A (baseline)          & 0.00 & 1.0 & $+0.005$ & 1.009 \\
B (mean shifted)      & 0.05 & 1.0 & $+0.045$ & 1.007 \\
C (volatility doubled)& 0.00 & 4.0 & $+0.008$ & 4.021 \\
\bottomrule
\end{tabular}
\end{table}

Comparisons: \textbf{A vs B} isolates the mean, \textbf{A vs C} isolates the
variance, \textbf{B vs C} combines both.

\begin{table}[H]
\centering
\caption{Stage 1 signature distances ($L^2$ between expected signatures) by
truncation level. The mean difference is present from level~1; the variance
difference switches on at level~2.}
\begin{tabular}{@{}lccccc@{}}
\toprule
Comparison & L1 & L2 & L3 & L4 & L5 \\
\midrule
A\_vs\_B (mean)     & 0.808 & 0.808 & 7.29 & 9.16 & 34.7 \\
A\_vs\_C (variance) & 0.054 & 31.25 & 31.36 & 860.3 & 917.8 \\
B\_vs\_C (combined) & 0.862 & 31.25 & 32.79 & 854.8 & 901.5 \\
\bottomrule
\end{tabular}
\end{table}

\begin{figure}[H]
\centering
\includegraphics[width=0.8\textwidth]{signature_distance_vs_level.png}
\caption{Stage 1: $L^2$ distance between expected signatures versus truncation
level. The mean comparison (A\_vs\_B) is already non-zero at level~1; the
variance comparison (A\_vs\_C) starts near zero and jumps at level~2.}
\end{figure}

\begin{figure}[H]
\centering
\includegraphics[width=0.9\textwidth]{sample_paths_A_B_C.png}
\caption{Example Stage 1 paths. Process~C (double volatility) visibly wanders
further; process~B drifts gently upward.}
\end{figure}

\paragraph{Reading.} The mean is a first-order (level-1) property, recovered
immediately; the variance is a second-order (level-2) property, appearing only
once level~2 is included. The continued growth of the raw distance at levels~4--5
is a \emph{scale effect} (higher-level coordinates are numerically larger), not
new distributional information --- a Gaussian carries no independent information
beyond its first two moments.

% ===========================================================================
\section{Stage 2: one-dimensional higher moments}
% ===========================================================================
Stage~2 matches the mean and variance across three processes so that levels~1--2
cannot separate them; any difference the signature detects must then come from a
\emph{higher} moment.

\begin{table}[H]
\centering
\caption{Stage 2 processes (4000 paths). Mean and variance are matched; only
skewness and kurtosis differ.}
\begin{tabular}{@{}lcccc@{}}
\toprule
Process & Mean & Variance & Skewness & Excess kurtosis \\
\midrule
G (Gaussian)             & $-0.00$ & 1.01 & $+0.01$ & $+0.03$ \\
S (skew-normal, $\alpha=8$) & $-0.00$ & 0.99 & $\mathbf{+0.92}$ & $+0.75$ \\
T (Student-$t$, $\mathrm{df}=5$) & $-0.00$ & 1.00 & $+0.05$ & $\mathbf{+4.44}$ \\
\bottomrule
\end{tabular}
\end{table}

\begin{table}[H]
\centering
\caption{Stage 2 signature distances by level. Skewness (G\_vs\_S) switches on at
level~3; kurtosis (G\_vs\_T) switches on at level~4.}
\begin{tabular}{@{}lccccc@{}}
\toprule
Comparison & L1 & L2 & L3 & L4 & L5 \\
\midrule
G\_vs\_S (skewness) & 0.011 & 0.119 & \textbf{1.781} & 6.089 & 21.83 \\
G\_vs\_T (kurtosis) & 0.033 & 0.033 & 0.749 & \textbf{6.006} & 9.914 \\
S\_vs\_T (combined) & 0.022 & 0.127 & 1.036 & 1.045 & 13.12 \\
\bottomrule
\end{tabular}
\end{table}

\begin{figure}[H]
\centering
\includegraphics[width=0.8\textwidth]{higher_moment_distance_vs_level.png}
\caption{Stage 2: distance versus level. The skew line (G\_vs\_S) switches on at
level~3; the heavy-tail line (G\_vs\_T) stays small through level~3 (the
Student-$t$ is symmetric) and switches on at level~4.}
\end{figure}

\begin{figure}[H]
\centering
\includegraphics[width=0.9\textwidth]{sample_paths_G_S_T.png}
\caption{Example Stage 2 paths for G, S and T.}
\end{figure}

\paragraph{Reading.} Each moment becomes visible at its own signature level:
mean~$\to$~1, variance~$\to$~2, skewness~$\to$~3, kurtosis~$\to$~4. To ``see''
the $k$-th moment one must truncate the signature at level~$\geq k$ --- a direct
demonstration of the supervisor's point that the truncation level matters and
that richer processes need deeper truncation.

% ===========================================================================
\section{Stage 3: high-dimensional Gaussian benchmark}
% ===========================================================================
\subsection{Motivation: from moments to geometry}
A one-dimensional path has no geometry: the signature essentially sees only the
endpoint, and higher levels behave like powers of a single number
(Eq.~\eqref{eq:1d}). Once a path has $d>1$ channels, level~2 onward contains
genuine cross-channel terms $S^{i,j}$ (and ordered higher-order interactions
$S^{i,j,k},\dots$), which encode \textbf{covariance, correlation blocks and
co-movement}. Stage~3 tests this and compares signatures against classical
statistical estimators.

\paragraph{Path meaning.} Each path is a multivariate return stream
$X_t = (r_t^1, \dots, r_t^d)$, where $r_t^i$ is the increment of \emph{synthetic
asset $i$}. Thus $d=20$ means 20 synthetic assets evolving together, and $d=10$
means 10.

\subsection{Configurations and processes}
Because the feature count grows as $d^{\,k}$, we use two complementary settings.
\begin{table}[H]
\centering
\caption{The two configurations and their signature feature dimensions
($d+d^2+\cdots+d^{\text{depth}}$, no constant term).}
\begin{tabular}{@{}llcccc@{}}
\toprule
Config & Purpose & $d$ & depth & Feature dim & Batch size \\
\midrule
\code{d20\_depth3} & wide, shallow: cross-channel geometry & 20 & 3 & $8{,}420$ & 250 \\
\code{d10\_depth5} & narrow, deep: higher-order interactions & 10 & 5 & $111{,}110$ & 100 \\
\bottomrule
\end{tabular}
\end{table}
All four processes are multivariate Gaussian random walks with increments
$\dd X \sim \mathcal{N}(\mu\,\dd t,\ \Sigma\,\dd t)$, $\dd t=1$.
\begin{table}[H]
\centering
\caption{The four high-dimensional processes.}
\renewcommand{\arraystretch}{1.2}
\begin{tabular}{@{}p{2.2cm}p{3.2cm}p{4.0cm}p{4.4cm}@{}}
\toprule
Process & Mean structure & Covariance structure & Interpretation \\
\midrule
A baseline & $0$ & identity & $d$ independent assets, no drift, unit vol \\
B mean-shift & first $25\%$ of channels $=0.05$ & identity & a subset of assets has a directional trend \\
C corr.-block & $0$ & identity $+$ correlated blocks & sector / block co-movement \\
D common-shock & $0$ & market factor $+$ idiosyncratic & a market-wide factor drives co-movement \\
\bottomrule
\end{tabular}
\end{table}
The correlation blocks in C set pairwise correlation $\rho=0.5$ on channels
$0$--$4$ and $\rho=0.3$ on channels $5$--$9$ (remaining channels independent).
Because only one property changes per process, each comparison isolates one
effect: A~vs~B (mean), A~vs~C and A~vs~D (covariance geometry), B~vs~C (mixed).

\begin{figure}[H]
\centering
\includegraphics[width=\textwidth]{sample_paths_first_5_dimensions.png}
\caption{Example high-dimensional paths, channels 0--4 (colour = channel), one
panel per process. Process~C's low channels visibly co-move (correlation blocks);
Process~D shows market-wide co-movement.}
\end{figure}

\subsection{Results: classical statistical distances}
\begin{table}[H]
\centering
\caption{Statistical distances between processes ($10{,}000$ paths, 300 steps).
The mean-shift (B) is seen only by the mean estimator; the correlation block~(C)
and common shock~(D) are seen only by the (increment) covariance estimator. The
increment covariance is the clean detector: the terminal covariance scales with
the number of steps and so has a large sampling-noise floor.}
\begin{tabular}{@{}llccc@{}}
\toprule
Config & Comparison & Mean-vec $L^2$ & Incr.\ cov.\ Frob. & Term.\ cov.\ Frob. \\
\midrule
\code{d20\_depth3} & A\_vs\_B (mean)  & \textbf{33.35} & 0.017 & 82.5 \\
\code{d20\_depth3} & A\_vs\_C (corr)  & 1.00 & \textbf{2.610} & 791.6 \\
\code{d20\_depth3} & A\_vs\_D (shock) & 1.09 & \textbf{1.798} & 532.7 \\
\midrule
\code{d10\_depth5} & A\_vs\_B (mean)  & \textbf{21.33} & 0.009 & 42.9 \\
\code{d10\_depth5} & A\_vs\_C (corr)  & 1.02 & \textbf{2.612} & 773.8 \\
\code{d10\_depth5} & A\_vs\_D (shock) & 0.88 & \textbf{0.901} & 280.5 \\
\bottomrule
\end{tabular}
\end{table}

\begin{figure}[H]
\centering
\includegraphics[width=\textwidth]{highdim_statistical_distances.png}
\caption{Classical statistical distances. The mean-shift comparisons have small
covariance distance; the correlation/shock comparisons have large covariance
distance and near-zero mean distance.}
\end{figure}

\subsection{Results: signature distances}
\begin{table}[H]
\centering
\caption{Standardised level-wise signature distances (each coordinate in units of
Process~A's variability). The mean shift (A\_vs\_B) peaks at level~1; the
covariance-geometry comparisons (A\_vs\_C, A\_vs\_D) are near-zero at level~1 and
peak at level~2.}
\begin{tabular}{@{}llccccc@{}}
\toprule
Config & Comparison & L1 & L2 & L3 & L4 & L5 \\
\midrule
\code{d20\_depth3} & A\_vs\_B & \textbf{0.433} & 0.135 & 0.084 & --- & --- \\
\code{d20\_depth3} & A\_vs\_C & 0.013 & \textbf{0.094} & 0.014 & --- & --- \\
\code{d20\_depth3} & A\_vs\_D & 0.014 & \textbf{0.064} & 0.015 & --- & --- \\
\midrule
\code{d10\_depth5} & A\_vs\_B & \textbf{0.388} & 0.109 & 0.098 & 0.033 & 0.028 \\
\code{d10\_depth5} & A\_vs\_C & 0.018 & \textbf{0.183} & 0.015 & 0.074 & 0.016 \\
\code{d10\_depth5} & A\_vs\_D & 0.016 & \textbf{0.068} & 0.015 & 0.028 & 0.016 \\
\bottomrule
\end{tabular}
\end{table}

\begin{figure}[H]
\centering
\includegraphics[width=\textwidth]{highdim_levelwise_signature_distance.png}
\caption{The key figure: standardised level-wise signature distance. A\_vs\_B
(mean shift) peaks at level~1; A\_vs\_C and A\_vs\_D (covariance geometry) are flat
at level~1 and peak at level~2. The mixed comparison B\_vs\_C is elevated at both.}
\end{figure}

\begin{figure}[H]
\centering
\includegraphics[width=\textwidth]{highdim_signature_distance_by_depth.png}
\caption{Raw whole-signature distance as depth grows (levels $1..k$ combined).
The raw distance is dominated by the magnitude of the highest included level ---
a scale effect --- which is exactly why the standardised view above is needed to
read the true structure.}
\end{figure}

\paragraph{The scale caveat, made concrete.} The raw and level-normalised
distances grow monotonically with level for \emph{every} comparison. For
instance, the raw A\_vs\_C level-3 distance in \code{d20\_depth3} is $\approx
2748$, dwarfing the level-2 value of $\approx 396$ --- purely a magnitude
artifact. Only after standardising by Process~A's coordinate standard deviation
does the genuine ``level-2 covariance'' signal emerge. Normalised or standardised
distances must therefore always be reported alongside the raw ones.

\subsection{Scaling}
A full run of both configurations at $100{,}000$ paths peaked at
$\approx 0.46$~GB resident memory --- flat regardless of path count, as the
streaming design intends --- and took $\approx 2$~min for \code{d20\_depth3} and
$\approx 26$~min for \code{d10\_depth5}. The extra paths only sharpen the results:
the A\_vs\_C level-1 noise floor falls from $0.013$ at $10{,}000$ paths to $0.004$
at $100{,}000$, making the level-2 covariance signal stand out even more clearly.

% ===========================================================================
\section{Stage 4: advanced financial processes}
% ===========================================================================
Stage~4 takes the six models of Section~\ref{sec:models} --- Black--Scholes plus
the four departures --- as 1-D log-price paths, all calibrated to the same baseline
volatility so that levels~1--2 are comparable and the interesting differences live
in the higher moments and the path memory. Time augmentation is on (effective
dimension~2), so the signature is sensitive to path ordering; at depth~5 the
feature vector has only $2+4+8+16+32=62$ coordinates.

\begin{table}[H]
\centering
\caption{Stage-4 processes and empirical statistics ($5{,}000$ paths, 100 steps,
horizon 1 year, seed 42). All share the baseline volatility; the differences are in
the higher moments and the path memory.}
\renewcommand{\arraystretch}{1.15}
\begin{tabular}{@{}llccc@{}}
\toprule
Code & Process (feature) & Term.\ var. & Incr.\ var. & Term.\ mean $\|\cdot\|$ \\
\midrule
\code{BS}    & Black--Scholes (baseline)        & 0.0401 & 0.000401 & 0.021 \\
\code{HEST}  & Heston (stochastic vol)          & 0.0421 & 0.000401 & 0.021 \\
\code{MERT}  & Merton (jumps)                   & 0.0718 & 0.000727 & 0.029 \\
\code{BATES} & Bates (jump-diffusion)           & 0.0736 & 0.000727 & 0.034 \\
\code{FBM}   & fractional BM ($H{=}0.70$, memory)& 0.0387 & \textbf{0.000063} & 0.002 \\
\code{RVOL}  & rough Bergomi ($H{=}0.10$, rough vol)& 0.0644 & 0.000427 & \textbf{0.297} \\
\bottomrule
\end{tabular}
\end{table}

\code{BS} attains its designed variance $\sigma^2 T = 0.04$. Jumps roughly double
the variance and shift the mean (negative-skew jumps). \code{FBM} has a much
\emph{smaller} per-step increment variance yet the \emph{same} terminal variance:
its positively correlated increments accumulate to the same endpoint spread --- the
defining fingerprint of long memory. \code{RVOL}'s large negative mean (the
$-\tfrac12 v\,\dd t$ It\^o term under a heavy-tailed variance) reflects rough,
spiky volatility.

\begin{table}[H]
\centering
\caption{Standardised level-wise signature distances from the Black--Scholes
baseline (each coordinate in units of \code{BS}'s variability). Bold marks where
each feature switches on: stochastic vol grows through the higher levels; jumps
switch on at level~2; Bates is the strongest of the diffusion family; long memory
(\code{FBM}) is modest but non-zero despite a matched terminal variance; rough
volatility dominates at every level.}
\begin{tabular}{@{}lccccc@{}}
\toprule
Comparison & L1 & L2 & L3 & L4 & L5 \\
\midrule
\code{BS\_vs\_HEST}  (stoch.\ vol) & 0.001 & 0.020 & 0.092 & 0.134 & \textbf{0.191} \\
\code{BS\_vs\_MERT}  (jumps)       & 0.040 & \textbf{0.276} & 0.373 & 0.449 & 0.545 \\
\code{BS\_vs\_BATES} (jump-diff.)  & 0.065 & \textbf{0.295} & 0.417 & 0.519 & 0.646 \\
\code{BS\_vs\_FBM}   (long memory) & \textbf{0.114} & 0.073 & 0.135 & 0.171 & 0.188 \\
\code{MERT\_vs\_BATES} (added vol) & 0.025 & 0.026 & 0.090 & 0.119 & 0.168 \\
\code{BS\_vs\_RVOL}  (rough vol)   & 1.367 & 1.266 & 1.587 & 1.991 & \textbf{2.608} \\
\code{FBM\_vs\_RVOL} (memory kinds)& 1.481 & 1.327 & 1.650 & 2.043 & \textbf{2.650} \\
\bottomrule
\end{tabular}
\end{table}

\begin{figure}[H]
\centering
\includegraphics[width=\textwidth]{advanced_levelwise_signature_distance.png}
\caption{Stage 4: standardised level-wise signature distance from Black--Scholes.
Rough volatility (\code{RVOL}) is the most extreme departure at every level; below
it, jumps (\code{MERT}, \code{BATES}) switch on at level~2 while stochastic vol
(\code{HEST}) and long memory (\code{FBM}) contribute smaller, higher-order
signals. (The axis label ``$\div$~Process-A coord std'' refers to the \code{BS}
baseline, reusing the Stage-3 plotting code.)}
\end{figure}

\begin{figure}[H]
\centering
\includegraphics[width=\textwidth]{advanced_sample_paths.png}
\caption{Example Stage-4 log-price paths, one panel per process. Jump processes
show discrete drops; the rough-vol path is visibly the most erratic.}
\end{figure}

\paragraph{The blind spot of moments, made concrete.} The classical
increment-covariance distance between \code{BS} and \code{FBM} is only
$\approx 3.4\times10^{-4}$ --- essentially the same as any other pair --- because
fractional Gaussian noise still has a well-defined marginal variance. The moments
are therefore nearly \emph{blind} to \code{FBM}'s long memory. Yet the signature
still separates \code{FBM} from \code{BS} (raw distance $0.029$, larger than the
stochastic-vol and jump pairs). This is the clearest single illustration in the
whole study of what signatures add over a moment/covariance table: they encode
\emph{ordered, path-dependent} structure that summary statistics cannot see.

% ===========================================================================
\section{Discussion}
% ===========================================================================
Across all four stages the signature behaves as the theory predicts, and in the
high-dimensional and advanced settings it agrees with the classical estimators
while adding structure a moment table alone cannot order:
\begin{itemize}
  \item \textbf{Mean shift} (A vs B) is a level-1 effect, recovered by both the
        mean estimator and signature level~1.
  \item \textbf{Covariance geometry} (A vs C correlation blocks; A vs D common
        shock) is a level-2+ effect carried by the cross terms $S^{i,j}$,
        recovered by both the covariance estimator and signature level~2.
  \item \textbf{Higher levels amplify scale}, so only normalised / standardised
        distances reveal the true switch-on structure.
  \item The signature additionally encodes \emph{ordered} interactions
        ($S^{i,j}\neq S^{j,i}$), a form of path geometry that symmetric
        covariance matrices cannot express.
  \item \textbf{Advanced processes} (Stage~4) confirm the pattern beyond the
        Gaussian world: stochastic volatility and jumps appear as higher-moment
        effects (jumps switching on at level~2), while \emph{path memory} ---
        long memory in \code{FBM} and roughness in \code{RVOL} --- is separated by
        the signature even when the marginal moments are matched and the classical
        increment-covariance estimator is blind to it.
\end{itemize}

% ===========================================================================
\section{Conclusion and next steps}
% ===========================================================================
Truncated path signatures characterise these processes by recovering their moment
and geometric information at interpretable levels: mean~$\to$~1, variance~$\to$~2,
skewness~$\to$~3, kurtosis~$\to$~4, and multi-channel covariance /
co-movement~$\to$~2+. Beyond the Gaussian world, they further separate stochastic
volatility, jumps and jump-diffusions, and --- most distinctively --- the
\emph{path memory} of fractional Brownian motion and rough volatility, which the
marginal moments cannot see. All experiments are fully streaming and scale to
$100{,}000$ paths at constant memory.

Deliberately out of scope for now (the natural next stages) are: applying the
pipeline to the real \code{SPY/QQQ/TLT} rolling windows; and signature kernels and
signature-based anomaly / regime detection.

\vspace{1em}
\noindent\textbf{Reproducibility.} All results use seed~42. The experiments are
run from the project root with:
\begin{quote}\ttfamily
python "Gaussian testing/run\_gaussian\_signature\_test.py"\\
python "Gaussian testing/run\_higher\_moment\_test.py"\\
python "HighDim testing/run\_highdim\_gaussian\_benchmark.py"\\
python "Advanced processes testing/run\_advanced\_processes\_test.py"
\end{quote}
Tables and figures are written under each folder's \code{results/} directory.

% ===========================================================================
\begin{thebibliography}{99}
% ===========================================================================
\bibitem{Chen1957} K.-T. Chen. Integration of paths, geometric invariants and a
  generalized Baker--Hausdorff formula. \emph{Annals of Mathematics},
  65(1):163--178, 1957.
\bibitem{Lyons1998} T. Lyons. Differential equations driven by rough signals.
  \emph{Revista Matem\'atica Iberoamericana}, 14(2):215--310, 1998.
\bibitem{LyonsCaruanaLevy2007} T. Lyons, M. Caruana, T. L\'evy. \emph{Differential
  Equations Driven by Rough Paths}. Lecture Notes in Mathematics 1908, Springer,
  2007.
\bibitem{FrizVictoir2010} P. Friz, N. Victoir. \emph{Multidimensional Stochastic
  Processes as Rough Paths}. Cambridge University Press, 2010.
\bibitem{HamblyLyons2010} B. Hambly, T. Lyons. Uniqueness for the signature of a
  path of bounded variation and the reduced path group. \emph{Annals of
  Mathematics}, 171(1):109--167, 2010.
\bibitem{BoedihardjoGengLyonsYang2016} H. Boedihardjo, X. Geng, T. Lyons, D. Yang.
  The signature of a rough path: uniqueness. \emph{Advances in Mathematics},
  293:720--737, 2016.
\bibitem{Fawcett2003} T. Fawcett. \emph{Problems in Stochastic Analysis:
  Connections between Rough Paths and Non-commutative Harmonic Analysis}. DPhil
  thesis, University of Oxford, 2003.
\bibitem{LevinLyonsNi2013} D. Levin, T. Lyons, H. Ni. Learning from the past,
  predicting the statistics for the future, learning an evolving system.
  arXiv:1309.0260, 2013.
\bibitem{ChevyrevKormilitzin2016} I. Chevyrev, A. Kormilitzin. A primer on the
  signature method in machine learning. arXiv:1603.03788, 2016.
\bibitem{ChevyrevOberhauser2022} I. Chevyrev, H. Oberhauser. Signature moments to
  characterize laws of stochastic processes. \emph{Journal of Machine Learning
  Research}, 23(176):1--42, 2022.
\bibitem{KidgerBonnier2019} P. Kidger, P. Bonnier, I. Perez Arribas, C. Salvi,
  T. Lyons. Deep signature transforms. \emph{Advances in Neural Information
  Processing Systems (NeurIPS)}, 2019.
\bibitem{KidgerLyons2021} P. Kidger, T. Lyons. Signatory: differentiable
  computations of the signature and logsignature transforms, on both CPU and GPU.
  \emph{International Conference on Learning Representations (ICLR)}, 2021.
\bibitem{Salvi2021} C. Salvi, T. Cass, J. Foster, T. Lyons, W. Yang. The signature
  kernel is the solution of a Goursat PDE. \emph{SIAM Journal on Mathematics of
  Data Science}, 3(3):873--899, 2021.
\bibitem{Heston1993} S. Heston. A closed-form solution for options with stochastic
  volatility with applications to bond and currency options. \emph{Review of
  Financial Studies}, 6(2):327--343, 1993.
\bibitem{CIR1985} J. Cox, J. Ingersoll, S. Ross. A theory of the term structure of
  interest rates. \emph{Econometrica}, 53(2):385--407, 1985.
\bibitem{Merton1976} R. Merton. Option pricing when underlying stock returns are
  discontinuous. \emph{Journal of Financial Economics}, 3(1--2):125--144, 1976.
\bibitem{Bates1996} D. Bates. Jumps and stochastic volatility: exchange rate
  processes implicit in Deutsche Mark options. \emph{Review of Financial Studies},
  9(1):69--107, 1996.
\bibitem{LordKoekkoekVanDijk2010} R. Lord, R. Koekkoek, D. van Dijk. A comparison
  of biased simulation schemes for stochastic volatility models. \emph{Quantitative
  Finance}, 10(2):177--194, 2010.
\bibitem{Andersen2008} L. Andersen. Simple and efficient simulation of the Heston
  stochastic volatility model. \emph{Journal of Computational Finance},
  11(3):1--42, 2008.
\bibitem{MandelbrotVanNess1968} B. Mandelbrot, J. Van Ness. Fractional Brownian
  motions, fractional noises and applications. \emph{SIAM Review},
  10(4):422--437, 1968.
\bibitem{GatheralJaissonRosenbaum2018} J. Gatheral, T. Jaisson, M. Rosenbaum.
  Volatility is rough. \emph{Quantitative Finance}, 18(6):933--949, 2018.
\bibitem{BayerFrizGatheral2016} C. Bayer, P. Friz, J. Gatheral. Pricing under rough
  volatility. \emph{Quantitative Finance}, 16(6):887--904, 2016.
\end{thebibliography}

\end{document}
